%% Finitely many eigenvalues determine the signature of a hyperbolic orbifold
%% Springer Nature LaTeX template v3.1 (December 2024), sn-jnl class, numbered Math and Physical
%% Sciences reference style, written to the conventions of Annals of Global Analysis and Geometry
%% (paper/jga/AGAG.md). No font package: text and mathematics in Computer Modern, as in the figures.
%% Figures: figures/out/E<k>.pdf, captions \figcapEOne ... in figures/captions.tex.
\documentclass[pdflatex,sn-mathphys-num]{sn-jnl}

\usepackage{graphicx}%
\usepackage{amsmath,amssymb,amsfonts}%
\usepackage{amsthm}%
\usepackage{mathtools}%
\usepackage{mathrsfs}%
\usepackage[title]{appendix}%
\usepackage{xcolor}%
\usepackage{booktabs}%
\usepackage{array}%
\usepackage{enumitem}%

\graphicspath{{../../}}% figures are addressed relative to the repository root

\theoremstyle{thmstyleone}%
\newtheorem{theorem}{Theorem}[section]%
\newtheorem{lemma}[theorem]{Lemma}%
\newtheorem{proposition}[theorem]{Proposition}%
\newtheorem{corollary}[theorem]{Corollary}%
\theoremstyle{thmstylethree}%
\newtheorem{definition}[theorem]{Definition}%
\newtheorem{remark}[theorem]{Remark}%
\newtheorem{problem}{Problem}%

\newcommand{\Orb}{\mathcal{O}}
\newcommand{\Cl}{\mathcal{C}}
\newcommand{\Sig}{\mathrm{Sig}}
\newcommand{\Area}{\operatorname{Area}}
\newcommand{\Hyp}{\operatorname{Hyp}}
\newcommand{\diam}{\operatorname{diam}}
\newcommand{\mult}{\operatorname{mult}}
\newcommand{\tr}{\operatorname{tr}}
\newcommand{\R}{\mathbb{R}}
\newcommand{\Z}{\mathbb{Z}}
\newcommand{\Q}{\mathbb{Q}}
\newcommand{\cZ}{\mathrm{Z}}
\newcommand{\Hb}{\mathbb{H}^2}

%% Figures: figures/out/<name>.pdf at the text width of 119 mm (figures/SPEC.md).
% Captions of the figures, one \caption-ready macro per figure, each at most two short sentences:
% what is shown, and the one point to take from it; the decoding a reader needs (line styles,
% markers, scales) is in the sentence of the text that cites the figure. Notation: theory/CONVENTIONS.md.
% Colours are named by tone only ("dark", "light") and only where needed to read the figure, so the
% captions hold in greyscale print. Needs \usepackage{amssymb} (\mathfrak). The \ref labels are those
% of paper/jga/manuscript.tex, which reads this file and attaches \figcapOne ... \figcapEight to
% F1-F8 (\sref is its macro for a label of the electronic supplement). F9 belongs to the companion
% note, paper/arith/note.tex, which carries the same caption in its own notation. Every statement in
% a caption is asserted by the figure's script in figures/src/ or proved in the paper. Old and new
% word counts: paper/jga/FIGURE-RESTORE.md.

% F1 (Introduction) -- figures/out/F1.pdf
\newcommand{\figcapOne}{One triangle of $\mathcal O(2,8,8)$ (a) and of $\mathcal O(3,3,12)$ (b), in schematic shapes, coloured by $4\pi t\,\mathfrak h_t(x,x)$ at $t=0.02$. The two orbifolds share their first two heat invariants, yet heat lingers differently near their cone points.}

% F2 (Section 5) -- figures/out/F2.pdf
\newcommand{\figcapTwo}{Exact tilings of the hyperboloid by the triangle groups of $\mathcal O(2,8,8)$ (a) and $\mathcal O(3,3,12)$ (b), in an oblique orthographic view. The coloured triangle and its pale mirror image form one copy of each orbifold.}

% F3 (Section 3) -- figures/out/F3.pdf
\newcommand{\figcapThree}{The mirror argument of Theorem~\ref{thm:sigsep} for $(2,8,8)$, $(3,3,12)$ (a) and $(1;15)$, $(0;3,3,5,5)$ (b). Discs mark $X^*$ and rings $-X^*$; they coincide only for an orbifold against itself.}

% F4 (Section 3) -- figures/out/F4.pdf
\newcommand{\figcapFour}{The largest $K_{\mathrm{mult}}$ on complete equal-area classes and on constructed pairs, against $s=\mathrm{Area}/2\pi$. The open growth question lies between the two step curves.}

% F5 (Section 7) -- figures/out/F5.pdf
\newcommand{\figcapFive}{(a) The difference of the heat traces of $\mathcal O(2,8,8)$ and $\mathcal O(3,3,12)$ from computed spectra, against its cone-point term. (b) Differences of heat traces in the family $(0;3,3,3,3)$, against the predicted geodesic term.}

% F6 (Section 4, what heat does not hear) -- figures/out/F6.pdf
\newcommand{\figcapSix}{(a) Four members of the family $(0;3,3,3,3)$ with their shortest closed geodesics; (b) its eigenvalues below $120$ at eight moduli. All members share every heat invariant, yet the spectrum moves.}

% F7 (Section 5) -- figures/out/F7.pdf
\newcommand{\figcapSeven}{Reciprocal sums $R$ of the hyperbolic triads of sum $S$, one band for each least order $p=2,\dots,8$. Adjacent bands first overlap at the dots and first collide at the circles.}

% F8 (Section 6) -- figures/out/F8.pdf
\newcommand{\figcapEight}{Recovery error of the cone orders against a worst-case error $\delta$ in every heat invariant, for $(2,3,7)$, $(2,8,8)$ and $(4,4,4)$. The slopes are the exponents of Theorem~\ref{thm:S3}.}

% F9 (companion note, Section 6) -- figures/out/F9.pdf; the note's own copy uses \N and \ciso.
\newcommand{\figcapNine}{(a) The real locus of $C_{27/2}$ in the chart $x+y+z=1$, with the permutations of
$(1,4,4)$ and $(1,1,4)$, that is, $(2,8,8)$ and $(3,3,12)$ up to scale. (b) The number $\mathcal N(S)$ of pairs of sum at most $S$ sharing $c_1,c_2$, divided by $S$
(heavy), and its fit (dashed), far above the isosceles lower bound, which grows like
$(c_{\mathrm{iso}}+o(1))\log S$ (thin).}

% E1-E3: figures of the eigen manuscript (paper/eigen/manuscript.tex), same conventions.
% E1 (Section 8) -- figures/out/E1.pdf
\newcommand{\figcapEOne}{(a) The eigenvalues $\lambda_1,\dots,\lambda_6$ of $\mathcal O(2,3,m)$ against $m$, with their upper bounds; (b) the same eigenvalues, rescaled as $h_m^2(\lambda_j-\frac14)/\pi^2$. As the cone point of order $m$ becomes cusp-like, the low spectrum crowds towards $\frac14$.}

% E2 (Section 6) -- figures/out/E2.pdf
\newcommand{\figcapETwo}{The gaps between the signature terms of the heat trace of $\mathcal O(2,8,8)$ and of its five competitors, with the error of the geodesic term and the error of keeping finitely many eigenvalues. Where half the nearest gap exceeds both errors, finitely many eigenvalues decide the signature.}

% E3 (Section 7) -- figures/out/E3.pdf
\newcommand{\figcapEThree}{Eigenvalues needed to decide the signature on computed spectra, observed and certified, against the a-priori count of Theorem~\ref{thm:E}, by systole. The certificate needs $10^2$ to $10^{11}$ times fewer eigenvalues.}

\newcommand{\paperfigure}[4][t]{%
  \begin{figure}[#1]
  \centering
  \includegraphics[width=119mm]{figures/out/#2.pdf}%
  \caption{#3}\label{#4}
  \end{figure}}

\raggedbottom
\hypersetup{pdftitle={Finitely many eigenvalues determine the signature of a hyperbolic orbifold},
  pdfauthor={Palaash Gang, Akshaj Agadi, Arjun Veluri, Jerry Wang, Jeremy Barreto, Aarin Chouthaiwale},
  pdfkeywords={hyperbolic orbifolds, Laplace eigenvalues, inverse spectral problems, heat trace, Selberg trace formula}}

\begin{document}

\title[Finitely many eigenvalues determine the signature]{Finitely many eigenvalues determine the signature of a hyperbolic orbifold}

\author*[1]{\fnm{Palaash} \sur{Gang}}\email{palaash.gang@indusschoolpune.com}
\author[2]{\fnm{Akshaj} \sur{Agadi}}\email{a0708@tamu.edu}
\author[1]{\fnm{Arjun} \sur{Veluri}}\email{arjun.v@indusschoolpune.com}
\author[3]{\fnm{Jerry} \sur{Wang}}\email{jerryiw@uci.edu}
\author[4]{\fnm{Jeremy} \sur{Barreto}}\email{jeremy.barreto@ltps.info}
\author[5]{\fnm{Aarin} \sur{Chouthaiwale}}\email{aarinchouthaiwale@gmail.com}

\affil*[1]{\orgname{Indus International School}, \orgaddress{\city{Pune}, \country{India}}}
\affil[2]{\orgname{Texas A\&M University}, \orgaddress{\city{College Station}, \state{Texas}, \country{USA}}}
\affil[3]{\orgname{University of California, Irvine}, \orgaddress{\city{Irvine}, \state{California}, \country{USA}}}
\affil[4]{\orgname{Lawrence High School}, \orgaddress{\city{Lawrenceville}, \state{New Jersey}, \country{USA}}}
\affil[5]{\orgname{DriveChange Learning and Resource Centre}, \orgaddress{\city{Pune}, \country{India}}}

\presentaddresstxt{ORCID iDs}
\presentaddress{P.~Gang 0009-0006-7448-9488; A.~Agadi 0009-0009-1378-0556; A.~Veluri 0009-0001-5940-3272; J.~Wang 0009-0004-1814-275X; J.~Barreto 0009-0006-6894-7323; A.~Chouthaiwale 0009-0008-3878-7183.}

\abstract{The Laplace spectrum of a closed orientable hyperbolic $2$-orbifold determines its signature, the genus and the orders of its cone points, but a measurement or a computation yields only finitely many eigenvalues, each to some accuracy. We show that this suffices, effectively: among orbifolds of area at most $A$, systole at least $\varepsilon$ and cone orders at most $M$, the first $N$ eigenvalues, each known to within $\delta$, determine the signature, with $N$ and $\delta$ given by explicit formulas. Compactness alone yields some $N$ and $\delta$; ours can be computed and checked, and an a-posteriori version certifies the signature of a computed spectrum from $21$ to $750$ eigenvalues in our examples, where the a-priori count ranges from $4\times10^4$ to $7\times10^{12}$. The proof turns the first difference of heat invariants, an integer multiple of an explicit rational number reached by the $M$-th invariant at the latest, into a gap between heat traces. The bound on the cone orders cannot be dropped: as $m$ grows, the orbifolds $\mathcal O(2,3,m)$ acquire arbitrarily many eigenvalues near $\frac14$, like a cusp.}

\keywords{Hyperbolic orbifolds, Laplace eigenvalues, Inverse spectral problems, Heat trace, Selberg trace formula, Effective bounds}

\pacs[MSC Classification]{58J53 (primary); 58J50, 35P15, 30F35, 57R18}

\maketitle

\section{Introduction}\label{sec:intro}

A drum is heard through its tones, and a closed hyperbolic surface or orbifold through the eigenvalues $0=\lambda_0<\lambda_1\le\lambda_2\le\cdots$ of its Laplacian. For a closed orientable hyperbolic $2$-orbifold $\Orb$ whose singular points are cone points, the whole spectrum determines the \emph{signature} $\sigma(\Orb)=(g;m_1,\dots,m_n)$, the genus of the underlying surface and the multiset of cone orders \cite[Thm~1.1, Prop.~3.3]{drydenstrohmaier2009}; isospectral orbifolds of the same signature need not be isometric \cite[Thm~A]{linowitzvoight2015}. An experiment or a computation, however, delivers finitely many eigenvalues, each with an error. This paper asks whether such data determine the signature, and with how many eigenvalues and what accuracy.

The answer is yes, with explicit numbers, once three quantities are bounded: the area, from below the systole (the least length of a closed geodesic), and from above the cone orders. Write $\Sig$ for the class of closed orientable hyperbolic $2$-orbifolds whose singular points are cone points, and, for $A,\varepsilon>0$ and an integer $M\ge2$,
\[
  \Cl(A,\varepsilon,M)=\{\Orb\in\Sig:\ \Area(\Orb)\le A,\ \ell(\Orb)\ge\varepsilon,\ \text{all cone orders at most }M\},
\]
with $\ell(\Orb)$ the systole; $\Sig(A,M)$ denotes the finite set of signatures of area at most $A$ with orders at most $M$.

\begin{theorem}[Finitely many eigenvalues]\label{thm:intro}
For $A>0$, $\varepsilon>0$ and an integer $M\ge2$, the numbers $N\ge1$ and $\delta>0$ of Theorem~\ref{thm:E}, explicit in $A$, $\varepsilon$ and $M$, have the following property. If $\Orb\in\Cl(A,\varepsilon,M)$ and $\tilde\lambda_0,\dots,\tilde\lambda_{N-1}$ satisfy $|\tilde\lambda_j-\lambda_j(\Orb)|\le\delta$, then $\sigma(\Orb)$ is the signature $\sigma\in\Sig(A,M)$ for which an explicit function $G_\sigma$ is nearest to $\sum_{j<N}e^{-\tilde\lambda_jt_*}$ at an explicit time $t_*$. In particular two orbifolds of $\Cl(A,\varepsilon,M)$ whose first $N$ eigenvalues agree to within $\delta$ have the same signature.
\end{theorem}

\emph{What effectivity adds.} Without constants the statement follows from compactness. Fix two signatures $\sigma\ne\sigma'$ in $\Sig(A,M)$. The orbifolds of signature $\sigma$ with systole at least $\varepsilon$ form a compact subset of moduli space, by Mumford's theorem \cite[Cor.~3]{mumford1971}, which Bers extended to Fuchsian groups with elliptic elements \cite{bers1972}; each $\lambda_j$ is continuous there, and two orbifolds of different signatures have different spectra \cite{drydenstrohmaier2009}. Covering the compact product of the two sets by the open sets where some $\lambda_j$ differ gives a finite $N$, then a $\delta>0$ by compactness again, and $\Sig(A,M)$ is finite. This argument gives no value of $N$ or $\delta$, and no way to check a given spectrum. Theorem~\ref{thm:E} gives both numbers as formulas. They are enormous (Table~\ref{tab:E}), but the same mechanism, run at one moderate time on the data instead of on worst-case bounds, gives a certificate that a computed spectrum can actually pass.

\begin{theorem}[A-posteriori certificate]\label{thm:intro-post}
Let $\Orb\in\Sig$, with systole at least $\ell$, diameter at most $\Delta$, area $A$ and signature in a finite set $\mathcal S$. Given approximations $\tilde\lambda_0\le\dots\le\tilde\lambda_N$ of its first $N+1$ eigenvalues with error bounds $\epsilon_j$, an explicit error budget $E_N(t)$ (Theorem~\ref{thm:post}) has the property: if, at some $t>0$, exactly one $\sigma\in\mathcal S$ satisfies $|\sum_{j<N}e^{-\tilde\lambda_jt}-G_\sigma(t)|\le E_N(t)$, then $\sigma(\Orb)=\sigma$.
\end{theorem}

On the computed spectra of $\Orb(2,8,8)$, $\Orb(3,3,12)$ and eight orbifolds of signature $(0;3,3,3,3)$ (Section~\ref{sec:practice}) the certificate succeeds with $21$ to $750$ eigenvalues, while Theorem~\ref{thm:E} asks for $4\times10^4$ to $7\times10^{12}$ (Table~\ref{tab:practice}).

The proof runs through the heat trace. By the Selberg trace formula the heat trace of $\Orb$ is a term $G_\sigma(t)$ fixed by the signature plus a positive contribution of the closed geodesics (Theorem~\ref{thm:trace}); the small-$t$ expansion of $G_\sigma$ has coefficients, the heat invariants, that are integer combinations of explicit rationals. At the first heat invariant where two signatures differ, the difference is an integer multiple of an explicit rational number (Lemma~\ref{lem:int}), and with orders at most $M$ that invariant is the $M$-th one at the latest (Theorem~\ref{thm:bounded}). Explicit, enveloping remainders for both expansions (Propositions~\ref{prop:remcone} and~\ref{prop:remarea}) turn this into a lower bound for $|G_\sigma(t)-G_{\sigma'}(t)|$ at a small explicit time; a diameter bound in terms of $A$, $\varepsilon$, $M$ (Theorem~\ref{thm:diam}) controls the closed geodesics there, and the counting function controls the eigenvalues beyond the $N$-th.

None of the three bounds is a matter of convenience for the method, and the bound on the orders is needed for the statement itself.

\begin{theorem}[High orders behave like cusps]\label{thm:intro-cusp}
The triangle orbifolds $\Orb(2,3,m)$, $m\ge7$, have area less than $\frac\pi3$ and systole at least $0.5620$, and $\lambda_j(\Orb(2,3,m))\le\frac14+\pi^2(j+1)^2/h_m^2$ with $h_m=\operatorname{arccosh}\frac1{2\sin(\pi/m)}\ge\log\frac m{2\pi}$. Consequently, for $A\ge\frac\pi3$ and $\varepsilon\le0.5620$, no $N$ and $\delta$ depending only on $A$ and $\varepsilon$ have the property of Theorem~\ref{thm:intro}.
\end{theorem}

The cone point of order $m$ sits at the tip of a cone of angle $2\pi/m$ and radius $h_m$, which for large $m$ looks like the end of a cusp, and the spectrum below any fixed level crowds towards $\frac14$, the bottom of the continuous spectrum of a cusp; Figure~\ref{fig:E1} shows the computed eigenvalues. Whether the systole bound is needed we do not know (Problem~\ref{prob:systole}).

\emph{Context.} The coefficients of the heat expansion of an orbifold are local \cite{donnelly1976}, \cite[Thm~4.8]{dggw2008}, and at constant curvature they are known in closed form \cite[Thm~4.20]{ucar2017}; Stanhope bounds the isotropy types of isospectral orbifolds \cite{stanhope2005}, and Dryden and Strohmaier show that the spectrum determines the signature \cite{drydenstrohmaier2009}. For Euclidean triangles, the first $N$ Dirichlet eigenvalues determine the triangle among triangles with all angles at least $\epsilon$, with $N$ depending on $\epsilon$ \cite{changdeturck1989}; the angle bound plays the part of our bounds on the systole and the orders. A companion manuscript in preparation, \emph{How much of a hyperbolic orbifold does heat hear?}, by the same authors, studies how many heat invariants determine the signature; it contains a sharper and more general form of Theorem~\ref{thm:bounded} below, which we prove here in the form we use. No proof here depends on it.

\emph{Organisation.} Section~\ref{sec:heat} derives the heat expansion with explicit remainders from the trace formula, and Section~\ref{sec:sig} the structure of the heat invariants of a signature. Section~\ref{sec:geom} bounds the diameter, Section~\ref{sec:count} counts eigenvalues, Section~\ref{sec:E} proves Theorem~\ref{thm:intro}, Section~\ref{sec:practice} the certificate, and Section~\ref{sec:hyp} discusses the hypotheses. Appendix~\ref{app:trace} extends the trace formula to the heat function.

\section{The heat trace and its expansion}\label{sec:heat}

\subsection{Signatures and the trace formula}

Let $\Orb\in\Sig$ have signature $\sigma=(g;m_1,\dots,m_n)$, $m_i\ge2$. It is hyperbolic exactly when $s(\sigma)=2g-2+\sum_i(1-\frac1{m_i})>0$, and then \cite[13.3.4--13.3.5]{thurston1980}, \cite[Thm~3.2]{drydenstrohmaier2009}
\begin{equation}\label{eq:area}
  \Area(\Orb)=2\pi s(\sigma)=2\pi\Big(2g-2+\sum_{i=1}^n\big(1-\tfrac1{m_i}\big)\Big).
\end{equation}
Write $\Orb=\Gamma\backslash\Hb$ with $\Gamma\subset PSL(2,\R)$ cocompact, so $\Gamma$ has no parabolic elements. Put $h_t(r)=e^{-t(1/4+r^2)}$ and $g_t(u)=e^{-t/4}e^{-u^2/4t}/\sqrt{4\pi t}$, so that $h_t(r)=\int g_t(u)e^{iru}du$. For hyperbolic $\gamma\in\Gamma$ let $\ell(\gamma)$ be its translation length and $\gamma=\gamma_0^k$ with $\gamma_0$ primitive; a \emph{closed geodesic} is a hyperbolic conjugacy class, and the systole $\ell(\Orb)$ is the least $\ell(\gamma)$. $\cZ_\Orb(t)=\sum_je^{-\lambda_jt}$ is the heat trace.

\begin{theorem}[Trace formula for the heat function]\label{thm:trace}
For every $t>0$, $\cZ_\Orb(t)=\mathrm I(t)+\sum_i\mathrm E_{m_i}(t)+\Hyp(t)$, all series and integrals converging absolutely, where
\begin{align*}
  \mathrm I(t)&=\frac{\Area(\Orb)}{4\pi}\int_\R r\tanh(\pi r)\,h_t(r)\,dr,\qquad
  \Hyp(t)=\sum_{[\gamma]}\frac{\ell(\gamma_0)\,g_t(\ell(\gamma))}{2\sinh(\ell(\gamma)/2)}>0,\\
  \mathrm E_m(t)&=\sum_{j=1}^{m-1}\frac1{2m\sin(\pi j/m)}\int_\R\frac{e^{-2\pi jr/m}\,h_t(r)}{1+e^{-2\pi r}}\,dr .
\end{align*}
\end{theorem}

\begin{proof}
For $h=\hat g$ with $g\in C_c^\infty(\R)$ even this is the Selberg trace formula for cocompact Fuchsian groups with elliptic elements as stated by Dryden and Strohmaier \cite[eq.~(1)]{drydenstrohmaier2009}, after \cite{hejhal1976,iwaniec2002}; Garbin and Jorgenson state it for the heat function \cite[Rem.~2.7, (2.8)]{garbinjorgenson2020}. Appendix~\ref{app:trace} gives the approximation argument, using Lemma~\ref{lem:counting}.
\end{proof}

For a signature $\sigma$ put $G_\sigma(t)=\mathrm I(t)+\sum_i\mathrm E_{m_i}(t)$, the part of $\cZ_\Orb$ fixed by $\sigma$.

\begin{lemma}[Counting closed geodesics]\label{lem:counting}
The number $n_\Orb(x)$ of closed geodesics of length at most $x$ satisfies $n_\Orb(x)\le2\pi(\cosh(x+3\diam)-1)/\Area(\Orb)\le\pi e^{x+3\diam}/\Area(\Orb)$, with $\diam$ the diameter of $\Orb$.
\end{lemma}

\begin{proof}
Let $F$ be the Dirichlet domain at a point $x_0$ fixed by no nontrivial element of $\Gamma$; $F\subset\bar B(x_0,\diam)$. Each hyperbolic class has a representative $\gamma$ whose axis meets $F$, at $y$ say, and $d(x_0,\gamma x_0)\le2d(x_0,y)+d(y,\gamma y)\le\ell(\gamma)+2\diam$. The translates $\gamma F$ with $d(x_0,\gamma x_0)\le x+2\diam$ have disjoint interiors and area $\Area(\Orb)$ each, and lie in $B(x_0,x+3\diam)$, of area $2\pi(\cosh(x+3\diam)-1)$.
\end{proof}

\begin{lemma}[The hyperbolic term]\label{lem:hyp}
Let $\ell$ be the systole. For $0<t\le\ell^2/(2(1+\ell))$,
\[
  0<\Hyp(t)\le B(\ell,\diam,t):=\frac{\pi e^{3\diam}\,\ell e^{\ell/2}}{\Area(\Orb)\,(1-e^{-\ell})}\Big(1+\frac{2t}{\ell-t}\Big)\frac{e^{-\ell^2/4t}}{\sqrt{4\pi t}},
\]
and on this range $B$ increases with $\diam$ and decreases with $\ell$. For every $t>0$,
\[
  0<\Hyp(t)\le\frac{2\sqrt\pi\,e^{3\diam}}{\Area(\Orb)\,(1-e^{-\ell})}\,e^{\frac{17t}4-\frac12-\ell}.
\]
\end{lemma}

\begin{proof}
With $\ell(\gamma_0)\le\ell(\gamma)$, $e^{-t/4}\le1$ and $2\sinh(x/2)\ge e^{x/2}(1-e^{-\ell})$ for $x\ge\ell$, $\Hyp(t)\le((1-e^{-\ell})\sqrt{4\pi t})^{-1}\int_{[\ell,\infty)}\phi\,dn_\Orb$ with $\phi(x)=xe^{-x/2-x^2/4t}$. On the range, $(\log\phi)'=1/x-1/2-x/2t<0$ for $x\ge\ell$, so integration by parts and Lemma~\ref{lem:counting} give $\int\phi\,dn_\Orb\le\frac{\pi e^{3\diam}}{\Area}\big(e^\ell\phi(\ell)+\int_\ell^\infty xe^{x/2-x^2/4t}dx\big)$, and as $x/(x-t)\le\ell/(\ell-t)$ on $[\ell,\infty)$ the integral is at most $\frac{2t\ell}{\ell-t}e^{\ell/2-\ell^2/4t}$. The logarithmic derivative of $B$ in $\ell$ is $-\frac12-(\frac\ell{2t}-\frac{1+\ell}\ell)-\frac{e^{-\ell}}{1-e^{-\ell}}-\frac{2t}{\ell^2-t^2}<0$, the bracket being nonnegative exactly on the range. For the second bound, $xe^{-x/2-x^2/4t}=(xe^{-x^2/8t})(e^{3x/2-x^2/8t})e^{-2x}\le2\sqrt t\,e^{-1/2}e^{9t/2}e^{-2x}$ for $x\ge0$ (maxima at $x=2\sqrt t$ and $x=6t$), and $\sum_{[\gamma]}e^{-2\ell(\gamma)}=2\int_\ell^\infty n_\Orb(x)e^{-2x}dx\le2\pi e^{3\diam-\ell}/\Area$.
\end{proof}

\subsection{The cone and area terms}

Write $\theta_j=\pi j/m$ and, for an integer $m\ge1$, $\Phi_m(u)=\sum_{j=1}^{m-1}(4m\sin\theta_j\sin(\theta_j-u))^{-1}$, analytic and even in $|u|<\pi/m$, $\Phi_1=0$.

\begin{lemma}[Closed form]\label{lem:Phi}
For $0<|u|<\pi/m$, $\Phi_m(u)=(\cot u-m\cot(mu))/(4m\sin u)$. With $u/\sin u=\sum_{i\ge0}\varsigma_iu^{2i}$ ($\varsigma_0=1$, $\varsigma_i>0$), $\Phi_m(u)=\sum_k\phi_k(m)u^{2k}$ with
\begin{equation}\label{eq:phik}
  m\,\phi_k(m)=\frac14\sum_{n=1}^{k+1}\varsigma_{k+1-n}\,\frac{4^n|B_{2n}|}{(2n)!}\,\big(m^{2n}-1\big),
\end{equation}
$B_{2n}$ the Bernoulli numbers. Every $\phi_k(m)$ is positive and strictly increasing in $m>1$, and $\phi_k(m)\le\Phi_m(\rho)\rho^{-2k}$ for $0<\rho<\pi/m$; at $\rho=\pi/2m$ this gives $\phi_k(m)\le\frac m4(2m/\pi)^{2k}$.
\end{lemma}

\begin{proof}
$(\sin\theta\sin(\theta-u))^{-1}=(\cot(\theta-u)-\cot\theta)/\sin u$, and the $\cot\theta_j$ cancel in pairs $j,m-j$. The function $\sum_{j=0}^{m-1}\cot(x+\pi j/m)-m\cot(mx)$ is $\pi/m$-periodic, has no poles and tends to $0$ as $\operatorname{Im}x\to\pm\infty$, so it vanishes; at $x=-u$ this is the closed form. From $\frac x2\coth\frac x2=\sum_nB_{2n}x^{2n}/(2n)!$ at $x=2iu$ and $\operatorname{sgn}B_{2n}=(-1)^{n+1}$, $\cot u-m\cot(mu)=\sum_{n\ge1}4^n|B_{2n}|(m^{2n}-1)u^{2n-1}/(2n)!$; multiplying by $u/\sin u=\prod_n(1-u^2/n^2\pi^2)^{-1}$ and by $1/(4mu)$ gives \eqref{eq:phik}, whose terms are positive and, as $(m^{2n}-1)/m$ is, strictly increasing in $m>1$. So $\Phi_m(\rho)=\sum_i\phi_i(m)\rho^{2i}$ is a sum of positive terms, each at most the sum; at $\rho=\pi/2m$, $\cot(m\rho)=0$ and $\Phi_m(\rho)=\cos\rho/(4m\sin^2\rho)\le m/4$ by $\sin\rho\ge2\rho/\pi$.
\end{proof}

Put, for $k\ge0$,
\[
  g_k(m)=(-1)^k\frac{(2k)!}{k!\,4^k}\phi_k(m),\quad b_l(m)=\sum_{k=0}^lg_k(m)\frac{(-1/4)^{l-k}}{(l-k)!},\quad\mu_k=4\int_0^\infty\frac{r^{2k+1}dr}{e^{2\pi r}+1},
\]
and let $\alpha_k$ be the coefficient of $t^{k-1}$ in $e^{-t/4}\big(t^{-1}-\sum_{j\ge0}(-1)^j\mu_jt^j/j!\big)$:
\[
  \alpha_k=\frac{(-1/4)^k}{k!}-\sum_{j=0}^{k-1}\frac{(-1)^j\mu_j}{j!}\,\frac{(-1/4)^{k-1-j}}{(k-1-j)!}.
\]
Every term of $b_l(m)$ has the sign $(-1)^l$, and every term of $\alpha_k$ the sign $(-1)^k$. (One has $\alpha_0,\dots,\alpha_3=1,-\frac13,\frac1{15},-\frac4{315}$; these are U\c{c}ar's smooth coefficients at curvature $-1$ \cite[(4.35)]{ucar2017}, and $b_0,b_1,b_2$ agree with the cone terms of \cite[Ex.~5.6]{dggw2008} and \cite[Thm~4.1]{schueth2019}; neither agreement is used.)

\begin{lemma}[Moments]\label{lem:moments}
For $0<a<2\pi$ let $F_a(r)=e^{-ar}/(1+e^{-2\pi r})$ and $\mathfrak m_n(a)=\int_\R r^nF_a(r)\,dr$. Then $\int_\R F_a(r)e^{sr}dr=(2\sin\frac{a-s}2)^{-1}$ for $a-2\pi<s<a$, and
\[
  \sum_{j=1}^{m-1}\frac{\mathfrak m_{2k}(2\theta_j)}{2m\sin\theta_j}=\frac{(2k)!}{4^k}\phi_k(m),\qquad \int_0^\infty\frac{4r^{2k+1}\,dr}{e^{2\pi r}+1}=\mu_k .
\]
\end{lemma}

\begin{proof}
With $u=e^{-2\pi r}$, $\int F_ae^{sr}dr=\frac1{2\pi}\int_0^\infty\frac{u^{(a-s)/2\pi-1}}{1+u}du=\frac1{2\sin((a-s)/2)}$ by Euler's beta integral. Differentiation under the integral at $s=0$ is dominated by $|r|^ne^{s_0|r|}F_a$ for $|s|\le s_0<\min(a,2\pi-a)$, so $\mathfrak m_n(a)=\partial_s^n(2\sin\frac{a-s}2)^{-1}|_{s=0}$, and the sum is $\partial_s^{2k}\Phi_m(\frac s2)|_{s=0}=\frac{(2k)!}{4^k}\phi_k(m)$. The second identity is the definition.
\end{proof}

\begin{proposition}[The cone term is enveloped]\label{prop:remcone}
For integers $m\ge2$, $K\ge0$ and every $t>0$, $\mathrm E_m(t)-\sum_{l<K}b_l(m)t^l$ has the sign $(-1)^K$ or vanishes, and $|\mathrm E_m(t)-\sum_{l<K}b_l(m)t^l|\le|b_K(m)|\,t^K$. Moreover $|b_K(m)|$ is strictly increasing in $m$.
\end{proposition}

\begin{proof}
$e^{t/4}\mathrm E_m(t)=\sum_j(2m\sin\theta_j)^{-1}\int F_{2\theta_j}(r)e^{-tr^2}dr$. For $x\ge0$ the Taylor remainder $e^{-x}-\sum_{k<K}(-x)^k/k!$ has the sign $(-1)^K$ and modulus at most $x^K/K!$; with $x=tr^2$, multiplying by the positive weight, integrating and summing, Lemma~\ref{lem:moments} shows that $R(t)=e^{t/4}\mathrm E_m(t)-\sum_{k<K}g_k(m)t^k$ has the sign $(-1)^K$ and $|R(t)|\le|g_K(m)|t^K$. Since
\[
  \mathrm E_m(t)-\sum_{l<K}b_lt^l=e^{-t/4}R(t)+\sum_{k<K}g_k(m)\,t^k\Big(e^{-t/4}-\sum_{i<K-k}\frac{(-t/4)^i}{i!}\Big),
\]
with $g_k$ of the sign $(-1)^k$ and the bracket of the sign $(-1)^{K-k}$ and modulus at most $(t/4)^{K-k}/(K-k)!$, every term has the sign $(-1)^K$ and the modulus of the sum is at most $t^K\sum_{k\le K}|g_k|4^{k-K}/(K-k)!=|b_K(m)|t^K$. Monotonicity follows from Lemma~\ref{lem:Phi}.
\end{proof}

\begin{proposition}[The area term is enveloped]\label{prop:remarea}
For integers $K\ge0$ and every $t>0$, $\frac{4\pi}{\Area}\mathrm I(t)-\sum_{k=0}^K\alpha_kt^{k-1}$ has the sign $(-1)^{K+1}$ or vanishes, and its modulus is at most $|\alpha_{K+1}|\,t^K$.
\end{proposition}

\begin{proof}
From $r\tanh\pi r=|r|-2|r|/(e^{2\pi|r|}+1)$, $\frac{4\pi}{\Area}\mathrm I(t)=e^{-t/4}(\frac1t-J(t))$ with $J(t)=4\int_0^\infty re^{-tr^2}(e^{2\pi r}+1)^{-1}dr$. As before, $J(t)-\sum_{k<K}\frac{(-1)^k\mu_k}{k!}t^k$ has the sign $(-1)^K$ and modulus at most $\frac{\mu_K}{K!}t^K$; $\frac1t(e^{-t/4}-\sum_{i\le K}\frac{(-t/4)^i}{i!})$ has the sign $(-1)^{K+1}$ and modulus at most $\frac{t^K}{4^{K+1}(K+1)!}$; and the cross terms $-\frac{(-1)^k\mu_k}{k!}t^k(e^{-t/4}-\sum_{i<K-k}\frac{(-t/4)^i}{i!})$ have the sign $(-1)^{K+1}$ and modulus at most $\frac{\mu_k4^{k-K}}{k!(K-k)!}t^K$. These moduli are those of the terms of $\alpha_{K+1}$, all of the sign $(-1)^{K+1}$, so they sum to $|\alpha_{K+1}|t^K$.
\end{proof}

So $\cZ_\Orb(t)\sim\sum_{j\ge1}c_j(\Orb)t^{j-2}$ as $t\downarrow0$, by Theorem~\ref{thm:trace} and Lemma~\ref{lem:hyp}, with the \emph{heat invariants}
\begin{equation}\label{eq:cj}
  c_1=\frac{\Area}{4\pi},\qquad c_j=\alpha_{j-1}\frac{\Area}{4\pi}+\sum_{i=1}^nb_{j-2}(m_i)\quad(j\ge2),
\end{equation}
functions $c_j(\sigma)$ of the signature; and, since all $b_K(m_i)$ have the sign $(-1)^K$, for $K\ge0$ and every $t>0$
\begin{equation}\label{eq:Grem}
  \Big|G_\sigma(t)-\sum_{j=1}^{K+1}c_j(\sigma)\,t^{j-2}\Big|\le t^K\Big(\frac{\Area}{4\pi}|\alpha_{K+1}|+\sum_i|b_K(m_i)|\Big).
\end{equation}

\begin{lemma}[Cone polynomials]\label{lem:conepoly}
For $l\ge0$, $p_l(m)=(-1)^lm\,b_l(m)=4^{-l}\sum_{k=0}^l\frac{(2k)!}{k!(l-k)!}m\phi_k(m)$ is an even polynomial with rational coefficients of degree exactly $2l+2$, with $p_l(1)=0$ and leading coefficient $a_l=|B_{2l+2}|/(2(l+1)!\,(2l+1))$.
\end{lemma}

\begin{proof}
The formula for $p_l$ is the definition of $b_l$. By \eqref{eq:phik}, $p_l$ is a combination with positive rational weights of $m^{2n}-1$, $1\le n\le l+1$; degree $2l+2$ occurs only for $k=l$, $n=l+1$, with coefficient $\frac1{4^l}\frac{(2l)!}{l!}\cdot\frac{4^l|B_{2l+2}|}{(2l+2)!}=a_l$.
\end{proof}

\section{Heat invariants of a signature}\label{sec:sig}

Write $R=\sum_i1/m_i$, $P_k=\sum_im_i^k$, $\psi_k(x)=x^{2k-1}-x^{-1}$ and $\Psi_k(m)=\sum_i\psi_k(m_i)$; $H_L(\sigma)=(c_1,\dots,c_L)$.

\begin{lemma}[Heat data of a signature]\label{lem:sigdata}
Let $\sigma=(g;m)$, $\sigma'=(g';m')$ be hyperbolic, with $n$ and $n'$ cone points, and $L\ge1$. Then $H_L(\sigma)=H_L(\sigma')$ if and only if $2g+n-R(m)=2g'+n'-R(m')$ and $\Psi_k(m)=\Psi_k(m')$ for $1\le k\le L-1$. In that case $d=R(m')-R(m)\in\Z$, and the paddings $U=m\uplus\{1\}^{\max(d,0)}$, $V=m'\uplus\{1\}^{\max(-d,0)}$ satisfy $R(U)=R(V)$, $P_j(U)=P_j(V)$ for odd $j\le2L-3$, and $|U|+|V|=2\max(n+g-g',n'+g'-g)$.
\end{lemma}

\begin{proof}
$c_1$ agrees exactly when the areas do, by \eqref{eq:area} the first equation. Write $p_l(x)/x=\sum_{k=1}^{l+1}a_{l,k}\psi_k(x)$, which holds with $a_{l,l+1}=a_l\ne0$ because $p_l$ is even with $p_l(1)=0$ (Lemma~\ref{lem:conepoly}); this triangular system shows that, at equal area, $c_2,\dots,c_L$ agree exactly when $\Psi_1,\dots,\Psi_{L-1}$ do. The rest is direct: $d=2(g'-g)+n'-n$, an entry $1$ changes neither $n-R$ nor any $\Psi_k$, and $\Psi_k=P_{2k-1}-R$.
\end{proof}

\begin{lemma}[Parity]\label{lem:parity}
In $\Q[x_1,\dots,x_N]$, for odd $j$ the elementary symmetric function $e_j$ lies in the ideal generated by the power sums $s_1,s_3,\dots,s_j$.
\end{lemma}

\begin{proof}
By Newton's identities $e_j$ is a rational combination of products $s_{i_1}\cdots s_{i_r}$ with $\sum i_t=j$, and for odd $j$ some $i_t$ is odd.
\end{proof}

\begin{theorem}[Separation]\label{thm:sep}
If two hyperbolic signatures of area $A$ share $c_1,\dots,c_L$ with $L\ge\lfloor A/\pi\rfloor+4$, they are equal.
\end{theorem}

\begin{proof}
Each cone point adds at least $\frac12$ to $s(\sigma)$, so $n+4g\le A/\pi+4$, and likewise for $\sigma'$; so $\kappa_0:=|U|+|V|\le2\lfloor A/\pi\rfloor+8\le2L$ in Lemma~\ref{lem:sigdata}. Remove the elements common to $U$ and $V$, leaving disjoint $U^*,V^*$ with the same equal sums, and let $X=U^*\uplus(-V^*)$, of size $\kappa\le\kappa_0$, even. If $X\ne\emptyset$, its odd power sums vanish up to $2L-3\ge\kappa-3$, so $e_j(X)=0$ for odd $j\le\kappa-3$ (Lemma~\ref{lem:parity}), and $e_{\kappa-1}(X)=e_\kappa(X)\sum_{x\in X}x^{-1}=0$; so $\prod_{x\in X}(z-x)$ is even, and $X=-X$, contradicting disjointness. So $U=V$ and, by Lemma~\ref{lem:sigdata}, $\sigma=\sigma'$.
\end{proof}

For an integer $M\ge2$ let $\Sig_{\le M}$ be the orbifolds in $\Sig$ with all cone orders at most $M$.

\begin{theorem}[Bounded cone orders]\label{thm:bounded}
Two hyperbolic signatures with all cone orders at most $M$ that share $c_1,\dots,c_M$ are equal.
\end{theorem}

\begin{proof}
Let $\nu(a)=\mult_m(a)-\mult_{m'}(a)$ for $2\le a\le M$ and $y_a=a^2$. By Lemma~\ref{lem:sigdata}, $2g+n-R(m)=2g'+n'-R(m')$ and $0=\Psi_k(m)-\Psi_k(m')=\sum_a\frac{\nu(a)}a(y_a^k-1)$ for $1\le k\le M-1$. Subtracting the column $a=1$ ($y_1=1$) from the others in the Vandermonde matrix $(y_a^k)_{0\le k\le M-1,\,1\le a\le M}$ shows $\det(y_a^k-1)_{1\le k\le M-1,\,2\le a\le M}=\prod_{1\le a<b\le M}(y_b-y_a)\ne0$. So $\nu=0$, and then $n=n'$, $R(m)=R(m')$ and $g=g'$.
\end{proof}

\begin{lemma}[Integrality at the first difference]\label{lem:int}
Let $\sigma\ne\sigma'$ be hyperbolic signatures with orders at most $M$, $k$ the least index with $c_k(\sigma)\ne c_k(\sigma')$, and $d_k=c_k(\sigma)-c_k(\sigma')$.
\begin{enumerate}[label=(\roman*)]
\item If the areas differ, $k=1$ and $|d_1|\ge1/(2L)$, where $L$ is the least common multiple of the orders of $\sigma$ and $\sigma'$.
\item If the areas are equal, to $A$ say, then $2\le k\le\min(\lfloor A/\pi\rfloor+4,M)$ and $d_k=(-1)^ka_{k-2}(P_{2k-3}(U)-P_{2k-3}(V))$, with $U,V$ the paddings of Lemma~\ref{lem:sigdata} for $L=k-1$; the difference of power sums is a nonzero integer, so $|d_k|\ge a_{k-2}$.
\end{enumerate}
\end{lemma}

\begin{proof}
(i) $c_1=\frac12(2g-2+n-R)$ and $R\in\frac1L\Z$. (ii) The bound on $k$ is Theorems~\ref{thm:sep} and~\ref{thm:bounded}. Put $l=k-2$. By Lemma~\ref{lem:sigdata} with $L=k-1$, $R(U)=R(V)$ and $\Psi_j(U)=\Psi_j(V)$ for $j\le l$, padding by $1$ changing no $C_l=\sum b_l$ as $p_l(1)=0$. The terms $\alpha_{k-1}A/4\pi$ cancel, and with the triangular expansion of the proof of Lemma~\ref{lem:sigdata}, $d_k=(-1)^l\sum_{j\le l+1}a_{l,j}(\Psi_j(U)-\Psi_j(V))=(-1)^la_l(P_{2l+1}(U)-P_{2l+1}(V))$, using $R(U)=R(V)$; this is nonzero as $d_k\ne0$.
\end{proof}

The bound $|d_k|\ge a_{k-2}$ is attained, for instance by $(0;4,4,4)$ and $(0;3,4,6)$, of equal area, with $d_2=-\frac1{12}$.

\section{The diameter}\label{sec:geom}

\subsection{Three facts of plane hyperbolic geometry}

We use the disc model, with $\cosh d(z,w)=1+2|z-w|^2/((1-|z|^2)(1-|w|^2))$, the upper half-plane, with $\cosh d(z,w)=1+|z-w|^2/(2\operatorname{Im}z\operatorname{Im}w)$, and the hyperboloid $\{x\in\R^{1,2}:\langle x,x\rangle=-1,\ x_0>0\}$ with $\langle x,y\rangle=-x_0y_0+x_1y_1+x_2y_2$, in which lines are the intersections with planes $n^\perp$, $n$ a unit spacelike normal \cite{beardon1983}.

\begin{lemma}\label{lem:H}
\leavevmode
\begin{enumerate}[label=(H\arabic*)]
\item A rotation by $\phi$ about $q$ moves a point at distance $r$ from $q$ by $d$ with $\sinh\frac d2=\sinh r\,|\sin\frac\phi2|$. A hyperbolic $\gamma$ moves a point at distance $r$ from its axis by $d$ with $\sinh\frac d2=\sinh\frac{\ell(\gamma)}2\cosh r$; in particular by at least $\ell(\gamma)$.
\item The product of the reflections in two distinct lines is a rotation by twice their angle about their intersection if they meet, a translation by twice their distance along their common perpendicular if they are ultraparallel, and has no fixed point in $\Hb$ if they are asymptotic.
\item Let $\lambda$ be the line through $p\ne q$, $d=d(p,q)$, $0<\alpha,\beta\le\frac\pi2$, and $L_1\ni p$, $L_3\ni q$ the lines making the angles $\alpha$, $\beta$ with the segment $[p,q]$, on the same side of $\lambda$. Put $X=\sin\alpha\sin\beta\cosh d-\cos\alpha\cos\beta$. If $X<1$, $L_1$ and $L_3$ meet, on that side, in a point $w$, and the triangle $pqw$ has angles $\alpha,\beta,\theta$ with $\cos\theta=X$; if $X>1$ they are ultraparallel at distance $h$ with $\cosh h=X$; if $X=1$ they are asymptotic.
\end{enumerate}
\end{lemma}

\begin{proof}
(H1) In the disc, rotate about $0$ a point $z$ with $|z|=\tanh\frac r2$: $|z-e^{i\phi}z|=2\tanh\frac r2|\sin\frac\phi2|$ and $1-|z|^2=\cosh^{-2}\frac r2$, so $\cosh d-1=8\sinh^2\frac r2\cosh^2\frac r2\sin^2\frac\phi2=2\sinh^2r\sin^2\frac\phi2$, and $\cosh d-1=2\sinh^2\frac d2$. In the half-plane let $\gamma(z)=e^Lz$, with axis $i\R_{>0}$; a point $z$ at distance $r$ from the axis has $|z|/\operatorname{Im}z=\cosh r$ (the distance to the axis is attained along the circle $|w|=|z|$), and $\cosh d(z,e^Lz)-1=(e^L-1)^2|z|^2/(2e^L(\operatorname{Im}z)^2)=2\sinh^2\frac L2\cosh^2r$.
(H2) Let $s_1,s_2$ be the reflections. If the lines meet at $p$ at the angle $\theta$, $s_2s_1$ fixes $p$, preserves orientation, and its differential at $p$ is the product of the linear reflections in two lines at the angle $\theta$, a rotation by $2\theta$; an isometry is determined by its value and differential at one point. If they are ultraparallel with common perpendicular $\mu$, meeting them at $q_1,q_2$, both reflections preserve $\mu$ and each side of it, and act on $\mu$ as the reflections in $q_1,q_2$; so $s_2s_1$ preserves the sides of $\mu$ and translates $\mu$ by $2d(q_1,q_2)$, which determines it as that translation. If they are asymptotic, $s_2s_1$ fixes their common ideal point; a fixed point in $\Hb$ would be fixed by $s_1$ and $s_2$ (it is the centre of a rotation that is a product of reflections in lines through it), hence would lie on both lines.
(H3) For two lines with unit normals $n,n'$: if they meet at $w$, then $n,n'\in T_w=w^\perp$, a Euclidean plane, are normal to the lines there, so $|\langle n,n'\rangle|=\cos\theta'<1$ with $\theta'$ the acute angle between them, and if they are oriented outwards from a convex region with a vertex at $w$ of interior angle $\theta$, $\langle n,n'\rangle=\cos(\pi-\theta)=-\cos\theta$. If they are ultraparallel with common perpendicular $\mu(s)=\cosh s\,x+\sinh s\,v$ meeting them at $s=0,h$, then $n,n'$ are $\pm\mu'(0)$, $\pm\mu'(h)$, and $|\langle\mu'(0),\mu'(h)\rangle|=\cosh h>1$. So the asymptotic case is $|\langle n,n'\rangle|=1$. Now take $p=(1,0,0)$, $\lambda=\{x_2=0\}$, $q=(\cosh d,\sinh d,0)$, the triangle on the side $x_2>0$, and $u_q=(\sinh d,\cosh d,0)$ the unit tangent of $\lambda$ at $q$. Outward unit normals (away from the side $x_2>0$ and from $q$, respectively $p$) of $\lambda$, $L_1$, $L_3$ are $-e_2$, $n_1=-\sin\alpha\,e_1+\cos\alpha\,e_2$, $n_3=\sin\beta\,u_q+\cos\beta\,e_2$; indeed $\langle-e_2,n_1\rangle=-\cos\alpha$ and $\langle-e_2,n_3\rangle=-\cos\beta$, as the rule above requires at $p$ and $q$. Then $\langle n_1,n_3\rangle=-\sin\alpha\sin\beta\cosh d+\cos\alpha\cos\beta=-X$, and $X\ge-\cos(\alpha+\beta)\ge-1$, with equality only if $\alpha=\beta=\frac\pi2$ and $d=0$. So if $X<1$ the lines meet; not on the side $x_2<0$, where the triangle would have angles $\pi-\alpha$ and $\pi-\beta$ summing to at least $\pi$; so on the side $x_2>0$, with $-\cos\theta=\langle n_1,n_3\rangle=-X$. The other two cases follow from $|\langle n_1,n_3\rangle|=|X|$.
\end{proof}

\subsection{The diameter bound}

A point $\tilde x\in\Hb$ with nontrivial stabiliser $\Gamma_{\tilde x}$ is \emph{elliptic}; $\Gamma_{\tilde x}$ is cyclic of order $m$, generated by the rotation by $2\pi/m$ about $\tilde x$, and its image in $\Orb$ is a cone point of order $m$.

\begin{lemma}[Cone points are separated]\label{lem:sep}
For $\Orb\in\Cl(A,\varepsilon,M)$, distinct elliptic points satisfy $d(\tilde p,\tilde q)\ge d_0(\varepsilon,M):=\min\{\frac\varepsilon2,\operatorname{arccosh}(1+\frac2{\pi^2M^2})\}$.
\end{lemma}

\begin{proof}
Let $a,b\le M$ be the orders of $\Gamma_{\tilde p},\Gamma_{\tilde q}$, $\alpha=\pi/a$, $\beta=\pi/b$, $d=d(\tilde p,\tilde q)$, and $\lambda,L_1,L_3,X$ as in (H3), with reflections $s_\lambda,s_1,s_3$. By (H2) $s_1s_\lambda$ and $s_\lambda s_3$ are rotations about $\tilde p$, $\tilde q$ by $2\alpha$, $2\beta$, which lie in the cyclic stabilisers (both senses do), so $\gamma=s_1s_3=(s_1s_\lambda)(s_\lambda s_3)\in\Gamma$. As $\Gamma$ has no parabolic element and $\gamma\ne1$, $X\ne1$ by (H2). If $X>1$, $\gamma$ is hyperbolic with $\ell(\gamma)=2h$ and $\cosh h=X\le\cosh d$, so $\varepsilon\le2d$. If $X<1$, $\gamma$ is a rotation by $2\theta$ about $w$, so $w$ is elliptic of some order $c\le M$ and $\theta=\pi k/c$, $1\le k<c$. By (H3)
\[
  \cosh d-1=\frac{\cos\theta+\cos(\alpha+\beta)}{\sin\alpha\sin\beta}=\frac{2\cos\frac{\theta+\alpha+\beta}2\cos\frac{\theta-\alpha-\beta}2}{\sin\alpha\sin\beta}.
\]
The defect $\pi-\alpha-\beta-\theta=\pi(abc-bc-ac-kab)/(abc)$ is positive with integer numerator, so at least $\pi/(abc)$, and $\cos\frac{\theta+\alpha+\beta}2\ge\frac1{abc}$ by $\sin x\ge\frac{2x}\pi$. As $\pi/M\le\theta$ and $\pi/M\le\alpha,\beta\le\frac\pi2$, $|\theta-\alpha-\beta|\le\pi-\frac\pi M$, so $\cos\frac{\theta-\alpha-\beta}2\ge\sin\frac\pi{2M}\ge\frac1M$; with $\sin\alpha\sin\beta\le\pi^2/(ab)$, $\cosh d-1\ge\frac2{\pi^2cM}\ge\frac2{\pi^2M^2}$.
\end{proof}

\begin{lemma}[Balls have area bounded below]\label{lem:balls}
For $\Orb\in\Cl(A,\varepsilon,M)$, $r_0=\frac12d_0(\varepsilon,M)$ and $\rho_1=\operatorname{arsinh}(\sinh r_0\sin\frac\pi M)$, every ball of radius $2r_0$ in $\Orb$ has area at least $v_0=2\pi\min\{(\cosh r_0-1)/M,\ \cosh\rho_1-1\}$.
\end{lemma}

\begin{proof}
If $d(\tilde y,\gamma\tilde y)\ge2r$ for every $\gamma\notin\Gamma_{\tilde y}$, the projection maps $\Gamma_{\tilde y}\backslash B(\tilde y,r)$ injectively onto the ball of radius $r$ about the image of $\tilde y$, of area $2\pi(\cosh r-1)/|\Gamma_{\tilde y}|$. Let $x\in\Orb$. If some cone point $p$ has $d(x,p)<r_0$, the ball of radius $2r_0$ about $x$ contains the ball of radius $r_0$ about $p$, and for $\gamma\notin\Gamma_{\tilde p}$, $\gamma\tilde p\ne\tilde p$ is elliptic, so $d(\tilde p,\gamma\tilde p)\ge2r_0$ (Lemma~\ref{lem:sep}). Otherwise $\Gamma_{\tilde x}=1$; a hyperbolic $\gamma$ moves $\tilde x$ by at least $\varepsilon\ge2\rho_1$ (H1), and an elliptic one is a rotation by $2\pi j/k$, $k\le M$, about a point at distance at least $r_0$, so moves $\tilde x$ by at least $2\operatorname{arsinh}(\sinh r_0\sin\frac\pi M)=2\rho_1$ (H1); and $\rho_1\le r_0$.
\end{proof}

\begin{theorem}[Diameter bound]\label{thm:diam}
Every $\Orb\in\Cl(A,\varepsilon,M)$ has $\diam\Orb<D(A,\varepsilon,M):=4r_0A/v_0\le\frac A\pi\max(4M,M^2)\max(\frac4\varepsilon,\frac{10M}3)$.
\end{theorem}

\begin{proof}
Let $\diam\Orb=d(x,y)=\Delta$ and lift to $\tilde x,\tilde y$ with $d(\tilde x,\tilde y)=\Delta$. The image $c$ of the segment, by arclength, satisfies $d(c(s),c(s'))=|s-s'|$: otherwise $d(x,y)<\Delta$. The points $c(4r_0i)$, $0\le i\le\lfloor\Delta/4r_0\rfloor$, are $4r_0$ apart, so the balls of radius $2r_0$ about them are disjoint, and Lemma~\ref{lem:balls} gives $(\lfloor\Delta/4r_0\rfloor+1)v_0\le A$. For the closed form, $\cosh r-1\ge r^2/2$, $\cosh\rho_1-1\ge(\cosh r_0-1)\sin^2\frac\pi M$ and $\sin\frac\pi M\ge\frac2M$ give $v_0\ge\pi r_0^2/\max(M,M^2/4)$, and $\operatorname{arccosh}(1+x)\ge\sqrt{2x/(1+x)}$ gives $d_0\ge\min(\frac\varepsilon2,\frac{0.6}M)$.
\end{proof}

The diameter enters only through Lemma~\ref{lem:counting}, the count of closed geodesics; in no row of Table~\ref{tab:E} is it the binding constraint through $t_1$, and it binds through $t_3$ only when $M$ is small. Table~\ref{tab:D} compares $D$ with true diameters. The bound grows like $A/\varepsilon$ and like $AM^3$; by Proposition~\ref{prop:233} below some growth in $M$ is necessary.

\begin{table}[h]
\centering\footnotesize
\caption{$D(A,\varepsilon,M)$ of Theorem~\ref{thm:diam}. The orbifolds $\Orb(2,8,8)$ and $\Orb(3,3,12)$ ($A=\pi/2$) have diameter below $4.9$ and $4.4$, the eight of signature $(0;3,3,3,3)$ of Section~\ref{sec:practice} ($A=4\pi/3$) at most $7.8$.}\label{tab:D}
\begin{tabular}{@{}lllll@{}}
\toprule
$A$ & $\varepsilon$ & $M=3$ & $M=12$ & $M=100$\\
\midrule
$\pi/2$ & $1$ & $56.6$ & $1125$ & $6.37\times10^5$\\
$4\pi/3$ & $1$ & $150.9$ & $3001$ & $1.70\times10^6$\\
$4\pi/3$ & $0.1$ & $640$ & $3184$ & $1.70\times10^6$\\
$10\pi$ & $1$ & $1132$ & $2.25\times10^4$ & $1.27\times10^7$\\
\bottomrule
\end{tabular}
\end{table}

\section{Counting eigenvalues}\label{sec:count}

\begin{lemma}[Elementary bounds]\label{lem:elem}
(i) Every $\Orb\in\Sig$ has $\Area(\Orb)\ge\pi/21$ and $n+4g\le\Area(\Orb)/\pi+4$. (ii) For $t>0$ and $m\ge2$, $0<\mathrm E_m(t)\le b_0(m)=\frac{m^2-1}{12m}$ and $0<\mathrm I(t)\le e^{-t/4}\Area/4\pi t$.
\end{lemma}

\begin{proof}
(i) With $s=\Area/2\pi$: $g\ge2$ gives $s\ge2$; $g=1$ needs $n\ge1$, $s\ge\frac12$; $g=0$, $n\ge5$ gives $s\ge\frac12$; $g=0$, $n=4$ has an order at least $3$, $s\ge\frac16$; and for $g=0$, $n=3$, $p\le q\le r$: $p\ge3$ forces $r\ge4$, $s\ge\frac1{12}$; $p=2$, $q=3$ forces $r\ge7$, $s\ge\frac1{42}$; $p=2$, $q=4$ forces $r\ge5$, $s\ge\frac1{20}$; $p=2$, $q\ge5$ gives $s\ge\frac1{10}$. The second statement holds as each cone point adds at least $\frac12$ to $s$. (ii) $0<h_t\le1$, and by Lemma~\ref{lem:moments} at $s=0$, $\mathrm E_m(t)\le\sum_j(4m\sin^2\theta_j)^{-1}=\Phi_m(0)=b_0(m)$; and $0\le r\tanh\pi r\le|r|$.
\end{proof}

\begin{proposition}[Counting and tails]\label{prop:counttail}
For every $\Orb$, $x>0$, $\Lambda>0$ and $0<s<t$: $\#\{j:\lambda_j\le x\}\le e\,\cZ_\Orb(1/x)$ and $\sum_{\lambda_j>\Lambda}e^{-\lambda_jt}\le e^{-\Lambda(t-s)}\cZ_\Orb(s)$. If $N\ge e\,\cZ_\Orb(1/\Lambda)$ is an integer, then $\lambda_j>\Lambda$ for $j\ge N$.
\end{proposition}

\begin{proof}
$\cZ(1/x)\ge\sum_{\lambda_j\le x}e^{-\lambda_j/x}\ge e^{-1}\#\{\lambda_j\le x\}$, and $e^{-\lambda t}\le e^{-\Lambda(t-s)}e^{-\lambda s}$ for $\lambda>\Lambda$.
\end{proof}

So, for $\Orb$ of area $A$, systole $\ge\ell$, diameter $\le\Delta$ and $n$ cone points of orders $m_i$, $\cZ_\Orb(s)\le\frac A{4\pi s}+\sum_ib_0(m_i)+\Hyp(s)$ for all $s$, with $\Hyp(s)$ bounded by Lemma~\ref{lem:hyp} with $\ell$ and $\Delta$ in place of the systole and diameter.

\section{Finitely many eigenvalues determine the signature}\label{sec:E}

Fix $A>0$, $\varepsilon>0$ and $M\ge2$. Put $n_*=\lfloor A/\pi\rfloor+4$, $k_*=\min(n_*,M)$, $L_M=\operatorname{lcm}(1,\dots,M)$, $\delta_1=1/(2L_M)$, $\delta_k=a_{k-2}$ for $2\le k\le k_*$, $\gamma_*=\min_{k\le k_*}\delta_k$,
\[
  Q(K)=\frac A{4\pi}|\alpha_{K+1}|+n_*\,|b_K(M)|,\qquad t_1=\min\Big\{1,\ \min_{1\le k\le k_*}\frac{\delta_k}{4\,Q(k-1)}\Big\},
\]
and $\Gamma(t)=\frac12\gamma_*t^{k_*-2}$.

\begin{lemma}[The gap]\label{lem:gap}
For $\sigma\ne\sigma'$ in $\Sig(A,M)$ and $0<t\le t_1$, $|G_\sigma(t)-G_{\sigma'}(t)|\ge\Gamma(t)$.
\end{lemma}

\begin{proof}
Let $k$ be the first index of difference, so $k\le k_*$ (Lemma~\ref{lem:int}). By \eqref{eq:Grem} with $K=k-1$, $n\le n_*$ (Lemma~\ref{lem:elem}) and $|b_K|$ increasing in $m$, $|G_\sigma(t)-\sum_{j\le k}c_j(\sigma)t^{j-2}|\le t^{k-1}Q(k-1)$, likewise for $\sigma'$, and so
\[
  |G_\sigma(t)-G_{\sigma'}(t)|\ge t^{k-2}\big(|d_k|-2tQ(k-1)\big)\ge\tfrac12t^{k-2}|d_k|\ge\tfrac12\delta_kt^{k-2}\ge\Gamma(t),
\]
since $|d_k|\ge\delta_k$ (Lemma~\ref{lem:int}, the lcm dividing $L_M$), $t\le\delta_k/(4Q(k-1))$, $t\le1$ and $k\le k_*$.
\end{proof}

\begin{theorem}[Finitely many eigenvalues determine the signature]\label{thm:E}
With $D=D(A,\varepsilon,M)$ of Theorem~\ref{thm:diam}, put
\begin{gather*}
  t_2=\frac{\varepsilon^2}{2(1+\varepsilon)},\qquad p=k_*-\tfrac32,\qquad\varpi=\frac{63\,\varepsilon e^{\varepsilon/2}}{(1-e^{-\varepsilon})\sqrt{4\pi}},\\
  y_0=2\Big(3D+\log\frac{16\varpi}{\gamma_*}+p\log\frac4{\varepsilon^2}+p\log\frac{2p}e\Big),\qquad t_3=\frac{\varepsilon^2}{4\max(y_0,1)},\qquad t_*=\min(t_1,t_2,t_3),\\
  \Gamma_*=\Gamma(t_*),\qquad\cZ^\sharp(s)=\frac A{4\pi s}+n_*\frac{M^2-1}{12M}+1,\qquad\Lambda=\frac2{t_*}\log\frac{8\cZ^\sharp(t_*/2)}{\Gamma_*},\\
  N=\big\lfloor e\,\cZ^\sharp(1/\Lambda)\big\rfloor+1,\qquad\delta=\min\Big\{\frac1{t_*},\frac{\Gamma_*}{8eNt_*}\Big\}.
\end{gather*}
Let $\Orb\in\Cl(A,\varepsilon,M)$ have eigenvalues $0=\lambda_0<\lambda_1\le\cdots$, counted with multiplicity and in increasing order, and let $\tilde\lambda_0,\dots,\tilde\lambda_{N-1}\ge0$ satisfy $|\tilde\lambda_j-\lambda_j|\le\delta$. Then $\sigma(\Orb)$ is the unique $\sigma\in\Sig(A,M)$ minimising $|\sum_{j<N}e^{-\tilde\lambda_jt_*}-G_\sigma(t_*)|$.
\end{theorem}

\begin{proof}
Write $\tilde\cZ=\sum_{j<N}e^{-\tilde\lambda_jt_*}$ and $\sigma_0=\sigma(\Orb)$. By Theorem~\ref{thm:trace},
$\tilde\cZ-G_{\sigma_0}(t_*)=\Hyp(t_*)-\sum_{j\ge N}e^{-\lambda_jt_*}+\sum_{j<N}(e^{-\tilde\lambda_jt_*}-e^{-\lambda_jt_*})$.
\emph{Hyperbolic term.} For $0<t\le t_2$, Lemma~\ref{lem:hyp} (with $\ell\ge\varepsilon$, $B$ decreasing in $\ell$ on the range), $\diam<D$, $\Area\ge\pi/21$ and $1+\frac{2t}{\varepsilon-t}\le3$ give $\Hyp(t)\le B_*(t):=\varpi e^{3D}t^{-1/2}e^{-\varepsilon^2/4t}$, which increases on $(0,t_2]$. With $y=\varepsilon^2/4t$, $B_*(t)\le\Gamma(t)/8$ reads $e^{-y}y^p\le\frac{\gamma_*}{16\varpi e^{3D}}(\varepsilon^2/4)^p$, and as $e^{-y/2}y^p\le(2p/e)^p$ it holds for $y\ge y_0$, so at $t=t_*$. Hence $\Hyp(s)\le\Gamma_*/8\le1$ for $s\le t_*$.
\emph{Tail.} For $s\le t_*$, $\cZ_\Orb(s)\le\cZ^\sharp(s)$ by Lemma~\ref{lem:elem} and the line above. As $\Lambda t_*\ge1$, $t_*/2$ and $1/\Lambda$ are at most $t_*$, so $N>e\,\cZ_\Orb(1/\Lambda)$, and Proposition~\ref{prop:counttail} with $s=t_*/2$ gives $\sum_{j\ge N}e^{-\lambda_jt_*}\le e^{-\Lambda t_*/2}\cZ^\sharp(t_*/2)=\Gamma_*/8$.
\emph{Perturbation.} $|e^{-at}-e^{-bt}|\le t|a-b|$ for $a,b\ge0$ gives at most $Nt_*\delta\le\Gamma_*/(8e)$.
So $|\tilde\cZ-G_{\sigma_0}(t_*)|<\frac38\Gamma_*$, while for $\sigma\ne\sigma_0$ Lemma~\ref{lem:gap} gives $|\tilde\cZ-G_\sigma(t_*)|>\Gamma_*-\frac38\Gamma_*$.
\end{proof}

Taking $\tilde\lambda_j=\lambda_j(\Orb')$ gives the last statement of Theorem~\ref{thm:intro}: orbifolds of $\Cl(A,\varepsilon,M)$ whose first $N$ eigenvalues agree to within $\delta$ have the same signature. The data must be the first $N$ eigenvalues, in order and with multiplicity; an approximation that misses one eigenvalue shifts every later index, and the theorem says nothing about it.

\begin{table}[h]
\centering\footnotesize
\caption{The constants of Theorem~\ref{thm:E}; the last column names the binding constraint on $t_*$.}\label{tab:E}
\setlength{\tabcolsep}{4pt}\begin{tabular}{@{}llllllllll@{}}
\toprule
$A$ & $\varepsilon$ & $M$ & $k_*$ & $D$ & $t_1$ & $t_3$ & $N$ & $\delta$ & binds\\
\midrule
$\pi/2$ & $1.8626$ & $8$ & $4$ & $402$ & $1.1\times10^{-7}$ & $3.5\times10^{-4}$ & $3.4\times10^{8}$ & $3.1\times10^{-21}$ & $t_1$\\
$\pi/2$ & $1.8626$ & $12$ & $4$ & $1125$ & $6.9\times10^{-9}$ & $1.3\times10^{-4}$ & $6.8\times10^{9}$ & $4.2\times10^{-25}$ & $t_1$\\
$4\pi/3$ & $2.634$ & $3$ & $3$ & $151$ & $9.3\times10^{-4}$ & $1.9\times10^{-3}$ & $4.3\times10^{4}$ & $1.5\times10^{-9}$ & $t_1$\\
$4\pi/3$ & $0.694$ & $3$ & $3$ & $151$ & $9.3\times10^{-4}$ & $1.3\times10^{-4}$ & $3.7\times10^{5}$ & $1.7\times10^{-10}$ & $t_3$\\
$4\pi/3$ & $0.694$ & $12$ & $5$ & $3001$ & $2.7\times10^{-11}$ & $6.7\times10^{-6}$ & $7.4\times10^{12}$ & $4.1\times10^{-41}$ & $t_1$\\
$2\pi$ & $0.1$ & $3$ & $3$ & $960$ & $7.8\times10^{-4}$ & $4.3\times10^{-7}$ & $2.4\times10^{8}$ & $2.7\times10^{-13}$ & $t_3$\\
$10\pi$ & $1$ & $3$ & $3$ & $1132$ & $3.3\times10^{-4}$ & $3.7\times10^{-5}$ & $1.1\times10^{7}$ & $5.6\times10^{-12}$ & $t_3$\\
$10\pi$ & $1$ & $12$ & $12$ & $2.25\times10^4$ & $2.5\times10^{-27}$ & $1.8\times10^{-6}$ & $3.8\times10^{30}$ & $8.7\times10^{-278}$ & $t_1$\\
\bottomrule
\end{tabular}
\end{table}

\emph{How large are $N$ and $\delta$.} Table~\ref{tab:E} gives them. Two constraints compete for $t_*$. One is the divergence of the cone expansion, through $t_1\approx\delta_k/Q(k-1)$, the least possible first difference against the largest possible next term; $|b_K(M)|$ grows roughly like $K!\,(M/\pi)^{2K}$, since $\Phi_M$ is singular at $\pm\pi/M$. The other is the geometry, through $t_3\approx\varepsilon^2/(24D)$, which keeps the closed geodesics below the gap. The bound $k_*\le M$ of Theorem~\ref{thm:bounded} matters: for $M=3$ it lowers $k_*$ from $\lfloor A/\pi\rfloor+4$ to $3$, and at $A=10\pi$ it changes $N$ from about $4\times10^{18}$ to $1.1\times10^7$. Once $t_*$ is fixed, $N\approx\frac{eA}{2\pi t_*}\log\frac{8\cZ^\sharp}{\Gamma_*}$ grows like $1/t_*$ and $\delta$ like $t_*^{k_*-2}$, up to logarithms; as $\varepsilon\to0$ with $A$, $M$ fixed, $t_3$ binds and $N$ grows by a factor near $8$ per halving of $\varepsilon$, roughly like $\varepsilon^{-3}\log\frac1\varepsilon$ (observations on the formulas, not proved asymptotics). We do not know the true order of the least admissible $N$.

Figure~\ref{fig:E2} shows the mechanism on the six signatures of area $\frac\pi2$ with orders at most $12$, for $\sigma_0=(0;2,8,8)$. The grey curves are $|G_{\sigma_0}(t)-G_\sigma(t)|$ for the five competitors, the darkest being $(0;3,3,12)$, whose difference starts at $d_3t=\frac{25}{12}t$; the dashed curve is the bound of Lemma~\ref{lem:hyp} on the closed geodesics, with the true systole and a diameter bound; and the dotted curves are the errors made by keeping $N=21$ and $N=100$ eigenvalues, the tail and perturbation terms of Theorem~\ref{thm:post}. Finitely many eigenvalues resolve the signatures at times where half the nearest gap exceeds the sum of both errors, the shaded window for $N=21$.

\paperfigure{E2}{\figcapETwo}{fig:E2}

\section{An a-posteriori certificate}\label{sec:practice}

\begin{theorem}[A-posteriori certificate]\label{thm:post}
Let $\Orb\in\Sig$ have area $A$, systole at least $\ell$, diameter at most $\Delta$, and signature in a finite set $\mathcal S$ of hyperbolic signatures of area $A$; put $\beta=\max_{\sigma\in\mathcal S}\sum_ib_0(m_i)$ and let $\mathcal H(s)$ be the bound of Lemma~\ref{lem:hyp} for $\Hyp(s)$, with $\ell$, $\Delta$ and the area $A$: the first bound $B(\ell,\Delta,s)$ for $s\le\ell^2/(2(1+\ell))$, where it is proved, and the bound for every $t>0$ otherwise. Let $\tilde\lambda_0\le\dots\le\tilde\lambda_N$ be nonnegative and $\epsilon_j\ge0$ with $|\tilde\lambda_j-\lambda_j(\Orb)|\le\epsilon_j$ for $j\le N$, the $\lambda_j$ counted with multiplicity in increasing order. For $t>0$ put
\[
  E_N(t)=\mathcal H(t)+\min_{0<s\le t}e^{-(\tilde\lambda_N-\epsilon_N)(t-s)}\Big(\frac A{4\pi s}+\beta+\mathcal H(s)\Big)+t\sum_{j<N}\epsilon_j .
\]
Then $|\sum_{j<N}e^{-\tilde\lambda_jt}-G_{\sigma(\Orb)}(t)|\le E_N(t)$ for every $t>0$. Hence if, at some $t$, exactly one $\sigma\in\mathcal S$ has $|\sum_{j<N}e^{-\tilde\lambda_jt}-G_\sigma(t)|\le E_N(t)$, in particular if $\sigma$ satisfies it and $E_N(t)$ is less than half the distance from $G_\sigma(t)$ to the nearest $G_{\sigma'}(t)$, $\sigma'\in\mathcal S\setminus\{\sigma\}$, then $\sigma(\Orb)=\sigma$.
\end{theorem}

\begin{proof}
As in the proof of Theorem~\ref{thm:E}: the hyperbolic term is at most $\mathcal H(t)$; the tail $\sum_{j\ge N}e^{-\lambda_jt}$ is at most $e^{-\lambda_N(t-s)}\cZ_\Orb(s)$, with $\lambda_N\ge\tilde\lambda_N-\epsilon_N$ and $\cZ_\Orb(s)\le\frac A{4\pi s}+\beta+\mathcal H(s)$ (Lemma~\ref{lem:elem}); and the perturbation is at most $t\sum\epsilon_j$. If $E_N(t)$ is less than half the nearest gap at $\sigma$ and $\sigma$ passes, every other $\sigma'$ is at distance more than $E_N(t)$.
\end{proof}

The hypotheses are a lower bound for the systole, an upper bound for the diameter, error bounds for the computed eigenvalues and the completeness of the computed list up to $\tilde\lambda_N$; for the triangle orbifolds below, twice the longest side bounds the diameter. We applied the certificate to the computed spectra of $\Orb(2,8,8)$ and $\Orb(3,3,12)$, about $2000$ eigenvalues each, complete below $1.6\times10^4$, and of eight orbifolds $\Orb_\vartheta$ of signature $(0;3,3,3,3)$, doubles of quadrilaterals with all angles $\frac\pi3$ and two mirror symmetries, with systoles from $2.634$ ($\vartheta=0$) down to $0.694$ ($\vartheta=2.8$) and $815$--$856$ eigenvalues each, complete up to $\lambda\approx2440$--$2560$. The spectra were computed with high-order finite elements and error estimates, and are in the public repository named in the data statement. The set $\mathcal S$ is all signatures of the same area with orders at most $M$. $N_{\rm apr}$ is the least $N$ for which the certificate succeeds at the best time on a grid from $0.002$ to $1$; $N_{\rm obs}$ is the least $N$ from which on, up to the end of the complete range, the data lie within half the nearest gap, the margin the certificate needs, judged against the data instead of against $E_N$.

\begin{table}[h]
\centering\footnotesize
\caption{Eigenvalues needed on computed spectra, against $N$ of Theorem~\ref{thm:E} for the same class; ranges are over the family members, and the $N_{\rm apr}$ range over the seven of systole at least $0.846$.}\label{tab:practice}
\begin{tabular}{@{}llllll@{}}
\toprule
orbifold & $M$ & $|\mathcal S|$ & $N_{\rm obs}$ & $N_{\rm apr}$ & $N$ of Theorem~\ref{thm:E}\\
\midrule
$\Orb(2,8,8)$ & $12$ & $6$ & $3$ & $21$ & $6.8\times10^9$\\
$\Orb(3,3,12)$ & $12$ & $6$ & $4$ & $39$ & $6.8\times10^9$\\
$\Orb_\vartheta$, $(0;3,3,3,3)$ & $3$ & $3$ & $3$--$50$ & $33$--$694$ & $4.3\times10^4$--$3.7\times10^5$\\
$\Orb_\vartheta$, $(0;3,3,3,3)$ & $12$ & $10$ & $5$--$98$ & $38$--$750$ & $7.4\times10^{12}$\\
\bottomrule
\end{tabular}
\end{table}

For the member of systole $0.694$ the certificate would need $t\le0.003$, where it needs more than the $815$ eigenvalues computed; this is a limit of the data, not of the method. Figure~\ref{fig:E3} plots the three counts against the systole of the member, for $M=3$ (open markers) and $M=12$ (filled), with the two triangle orbifolds at their own systoles in their colours: $N_{\rm obs}$ (circles), $N_{\rm apr}$ (diamonds) and the $N$ of Theorem~\ref{thm:E} (squares). The certificate needs between $10^2$ and $10^{11}$ times fewer eigenvalues than the a-priori count, and both practical counts grow as the systole falls, since closed geodesics then enter the heat trace earlier.

\paperfigure{E3}{\figcapEThree}{fig:E3}

\section{The hypotheses}\label{sec:hyp}

\subsection{The orders: high orders behave like cusps}

We need the case of J{\o}rgensen's inequality \cite{jorgensen1976} in which one element is hyperbolic, with his argument; we could not obtain the text of \cite{jorgensen1976} and prove what we use.

\begin{lemma}[J{\o}rgensen's inequality, hyperbolic case]\label{lem:jorgensen}
Let $\Gamma\subset PSL(2,\R)$ be discrete with no parabolic elements, $A\in\Gamma$ hyperbolic and $B\in\Gamma$ an element that fixes neither fixed point of $A$ in $\partial\Hb$ and does not interchange them. Then, for any lifts to $SL(2,\R)$, $|\tr^2A-4|+|\tr[A,B]-2|\ge1$.
\end{lemma}

\begin{proof}
Conjugate so that $A=\operatorname{diag}(u,u^{-1})$, $u>1$, with fixed points $0,\infty$, and let $B=\begin{psmallmatrix}a&b\\c&d\end{psmallmatrix}$. A direct computation gives $\tr[A,B]-2=-bc(u-u^{-1})^2$, so the claim is $\mu:=(u-u^{-1})^2(1+|bc|)\ge1$; suppose $\mu<1$. Let $B_0=B$, $B_{n+1}=B_nAB_n^{-1}=\begin{psmallmatrix}a_{n+1}&b_{n+1}\\c_{n+1}&d_{n+1}\end{psmallmatrix}$ and $x_n=b_nc_n$. Then
\[
  B_{n+1}=\begin{pmatrix}ua_nd_n-u^{-1}b_nc_n&-a_nb_n(u-u^{-1})\\c_nd_n(u-u^{-1})&-ub_nc_n+u^{-1}a_nd_n\end{pmatrix},\qquad x_{n+1}=-x_n(1+x_n)(u-u^{-1})^2 .
\]
So $|x_n|\le|x_0|$ implies $|x_{n+1}|\le\mu|x_n|$, and by induction $|x_n|\le\mu^n|x_0|\to0$. No $x_n$ vanishes: $x_0\ne0$ as $B$ fixes neither $0$ nor $\infty$; if $n\ge1$ is least with $x_n=0$, then $x_{n-1}=-1$, so $a_{n-1}d_{n-1}=0$; if both vanish, $B_{n-1}$ interchanges $0$ and $\infty$, which is excluded for $n-1=0$ and impossible for the hyperbolic $B_{n-1}$, $n-1\ge1$; if exactly one vanishes, $B_n$, a hyperbolic element with fixed points $B_{n-1}(0)$ and $B_{n-1}(\infty)$, shares exactly one fixed point with $A$, and then the commutator $[A,B_n]$ is a nontrivial parabolic element of $\Gamma$ (upper or lower unitriangular with a nonzero entry). Now $B_n$, $n\ge1$, is hyperbolic with repelling fixed point $p_n=B_{n-1}(0)=b_{n-1}/d_{n-1}$ and attracting fixed point $q_n=B_{n-1}(\infty)=a_{n-1}/c_{n-1}$, both finite and nonzero, with $p_n/q_n=x_{n-1}/(1+x_{n-1})\to0$, and with translation length that of $A$. Choose $k_n\in\Z$ with $u^{-2k_n}\in[\rho_n,u^2\rho_n)$, $\rho_n=|p_nq_n|^{-1/2}$, and put $C_n=A^{-k_n}B_nA^{k_n}\in\Gamma$; its fixed points $u^{-2k_n}p_n$ and $u^{-2k_n}q_n$ have moduli at most $u^2|p_n/q_n|^{1/2}\to0$ and at least $|q_n/p_n|^{1/2}\to\infty$. So $C_n\to A$ in $PSL(2,\R)$, while $C_n\ne A$ as $C_n$ does not fix $0$. This contradicts discreteness.
\end{proof}

For a rotation $\beta$ by $\phi$ about a point at distance $r$ from the axis of a hyperbolic $\gamma$ of translation length $L$, $\tr^2\gamma-4=4\sinh^2\frac L2$ and $\tr[\gamma,\beta]-2=4\sinh^2\frac L2\sin^2\frac\phi2\cosh^2r$ for any lifts: with $\gamma=\operatorname{diag}(e^{L/2},e^{-L/2})$ and $\beta=hkh^{-1}$, $k=\begin{psmallmatrix}\cos\frac\phi2&\sin\frac\phi2\\-\sin\frac\phi2&\cos\frac\phi2\end{psmallmatrix}$, $h=\begin{psmallmatrix}\cosh\frac r2&\sinh\frac r2\\\sinh\frac r2&\cosh\frac r2\end{psmallmatrix}$ (which moves $i$ along the unit circle, perpendicular to the axis, by $r$), the off-diagonal entries of $\beta$ are $\pm\sin\frac\phi2\cosh r$, so $bc=-\sin^2\frac\phi2\cosh^2r$, and $\tr[\gamma,\beta]-2=-bc(e^{L/2}-e^{-L/2})^2$.

\begin{proposition}[The family $\Orb(2,3,m)$]\label{prop:233}
For $m\ge7$ put $s_\infty=\operatorname{arccosh}\frac2{\sqrt3}$, $h_m=\operatorname{arccosh}\frac1{2\sin(\pi/m)}$ and $\sigma_*=2\operatorname{arsinh}\big(2\sqrt{1+\frac34\cosh^2(2s_\infty)}\big)^{-1}=0.56206\ldots$. Then $\Area\Orb(2,3,m)=2\pi(\frac16-\frac1m)<\frac\pi3$, $\ell(\Orb(2,3,m))\ge\sigma_*$, $\diam\Orb(2,3,m)\ge h_m\ge\log\frac m{2\pi}$, and the open ball of radius $h_m$ about the cone point of order $m$ is a hyperbolic cone of angle $2\pi/m$. So $\Orb(2,3,m)\in\Cl(\frac\pi3,\sigma_*,m)$.
\end{proposition}

\begin{proof}
$\Orb=\Orb(2,3,m)$ is the double of the triangle $T$ with vertices $v_2,v_3,v_m$ and angles $\frac\pi2,\frac\pi3,\frac\pi m$; its group $\Gamma_m$ is the orientation-preserving subgroup of the group $\Gamma^*$ generated by the reflections in the sides of $T$, of which $T$ is a fundamental domain \cite{beardon1983}, and the folding $f:\Hb\to T$ is $1$-Lipschitz with $d(z,\Gamma^*v)=d(f(z),v)$ for each vertex $v$. By (H3), $\cosh d(v_2,v_3)=2\cos(\pi/m)/\sqrt3$, so $s_m=d(v_2,v_3)<s_\infty$, and $\cosh d(v_2,v_m)=1/(2\sin(\pi/m))$, so $d(v_2,v_m)=h_m$; the folding of $\Orb$ onto $T$ is $1$-Lipschitz and the identity on $T$, so the cone points of orders $2$ and $m$ are at distance $h_m$ in $\Orb$, and $h_m\ge\log\frac1{2\sin(\pi/m)}\ge\log\frac m{2\pi}$.
\emph{The cone ball.} Let $P$ be the union of the $2m$ tiles $gT$, $g\in\Gamma^*$, with vertex $v_m$, the translates under the stabiliser of $v_m$ in $\Gamma_m$ of the fundamental domain $T\cup s(T)$, $s$ the reflection in $[v_m,v_2]$. For $\gamma\in\Gamma_m$ outside that stabiliser, $\gamma P$ and $P$ have disjoint interiors. $P$ is star-shaped about $v_m$, and its boundary consists of images of $[v_2,v_3]$, at distance $h_m$ from $v_m$ because the angle of $T$ at $v_2$ is $\frac\pi2$. So $B(v_m,h_m)\subset\operatorname{int}P$, its translates are equal or disjoint, and its image is the cone of angle $2\pi/m$ and radius $h_m$.
\emph{The systole.} Let $\gamma\in\Gamma_m$ be hyperbolic with axis $\alpha$ and $L=\ell(\gamma)$. As $\alpha$ is connected and unbounded, it is not covered by the disjoint open balls $gB(v_m,h_m)$; take $\tilde x\in\alpha$ outside them and $x=f(\tilde x)$, so $d(x,v_m)\ge h_m$. Write $x\in[v_m,q]$ with $q\in[v_2,v_3]$; as $d(v_m,\cdot)$ is convex on $[v_2,v_3]$ with minimum $h_m$ at $v_2$, $d(x,q)\le d(v_m,v_3)-h_m\le s_m$, so $d(x,v_3)\le2s_m$. So some elliptic point $\tilde w\in\Gamma^*v_3$, fixed by a rotation $\beta\in\Gamma_m$ by $\frac{2\pi}3$, lies at distance $r\le2s_m<2s_\infty$ from $\alpha$. The rotation $\beta$ fixes no point of $\partial\Hb$ and cannot interchange the endpoints of $\alpha$ (then $\beta^2$, a rotation of order $3$, would fix them), so Lemma~\ref{lem:jorgensen} and the trace identities give $4\sinh^2\frac L2(1+\frac34\cosh^2r)\ge1$, hence $L\ge\sigma_*$.
\end{proof}

On a region with metric $d\rho^2+w(\rho)^2d\vartheta^2$ and $w''=w$, the function $f=w^{-1/2}u(\rho)$ satisfies
\begin{equation}\label{eq:rayleigh}
  f'^2w=u'^2+q\,u^2+\frac d{d\rho}\Big(-\frac{w'}{2w}u^2\Big),\qquad q=\frac12-\frac{w'^2}{4w^2}=\begin{cases}\frac14-\frac1{4\sinh^2\rho},&w=\sinh\rho,\\\frac14+\frac1{4\cosh^2\rho},&w=\cosh\rho,\end{cases}
\end{equation}
as $f'^2w=u'^2-\frac{w'}wuu'+\frac{w'^2}{4w^2}u^2$ and $\frac d{d\rho}(-\frac{w'}{2w}u^2)=-(\frac{w''}{2w}-\frac{w'^2}{2w^2})u^2-\frac{w'}wuu'$.

\begin{proposition}[Unbounded orders]\label{prop:N1}
For $m\ge7$ and $j\ge1$, $\lambda_j(\Orb(2,3,m))\le\frac14+\pi^2(j+1)^2/h_m^2$. Moreover, for every $N\ge1$ and $\delta>0$, among any $K+1$ distinct orders $m\ge7$, where $K=(\lfloor\Lambda_N/\delta\rfloor+1)^{N-1}$ and $\Lambda_N=\frac14+\pi^2N^2/h_7^2$, there are two, $m\ne m'$, with $|\lambda_j(\Orb(2,3,m))-\lambda_j(\Orb(2,3,m'))|<\delta$ for all $j<N$.
\end{proposition}

\begin{proof}
On the cone ball of Proposition~\ref{prop:233}, with metric $d\rho^2+\sinh^2\rho\,d\vartheta^2$, $\vartheta\in\R/\frac{2\pi}m\Z$, cut $[a,h_m]$, $0<a<h_m$, into $j+1$ intervals of length $\ell=(h_m-a)/(j+1)$ and let $f_i=(\sinh\rho)^{-1/2}u_i$ with $u_i$ a sine arch on the $i$-th interval, zero elsewhere. These Lipschitz functions have disjoint supports, and by \eqref{eq:rayleigh}, as $u_i$ vanishes at the ends and $q\le\frac14$, $\int|\nabla f_i|^2/\int f_i^2\le\frac14+\pi^2/\ell^2$. The min--max principle on their span gives the bound; let $a\to0$. As $h_m$ increases with $m$, $\lambda_j\in[0,\Lambda_N]$ for $j<N$, and two of $K+1$ vectors $(\lambda_1,\dots,\lambda_{N-1})$ lie in a common product of the $\lfloor\Lambda_N/\delta\rfloor+1$ intervals $[i\delta,(i+1)\delta)$.
\end{proof}

\begin{remark}[The order bound cannot be dropped]\label{rem:Mneeded}
Let $A\ge\frac\pi3$ and $\varepsilon\le\sigma_*$. For every $N$ and $\delta$ there is $M$, for instance $M=(\lfloor\Lambda_N/\delta\rfloor+1)^{N-1}+7$, such that $\Cl(A,\varepsilon,M)$ contains $\Orb(2,3,m)$ and $\Orb(2,3,m')$, of different signatures, whose first $N$ eigenvalues agree to within $\delta$; whatever rule decides the signature from such data is wrong for one of them. So $N$ and $\delta$ in Theorem~\ref{thm:E} cannot depend on $A$ and $\varepsilon$ alone; an area bound and a systole bound do not bound the orders. For $A<\frac\pi3$ the question does not arise, since then $g=0$, $n=3$ and, by the case analysis of Lemma~\ref{lem:elem}, only finitely many triples occur.
\end{remark}

Figure~\ref{fig:E1} shows the computed eigenvalues $\lambda_1,\dots,\lambda_6$ of $\Orb(2,3,m)$ for $21$ orders from $7$ to $4096$, on logarithmic axes, with the bounds of Proposition~\ref{prop:N1} (thin dashed, one for each $j$) and the line $\frac14$. The spectrum of $\Orb(2,3,m)$ is the union of the Neumann and Dirichlet spectra of the triangle; we computed both by high-order finite elements at two mesh levels, which agree to at least $2.3\times10^{-10}$ relative. Every eigenvalue lies below its bound and decreases in $m$, and $\lambda_1$ falls from $44.89$ at $m=7$ to $0.473$ at $m=4096$; the computed values suggest $\lambda_j\to\frac14$, which Proposition~\ref{prop:N1} does not prove, and the plotted quantity $h_m^2(\lambda_j-\frac14)/\pi^2$ in panel (b) approaches $j^2$ slowly, the behaviour of an interval of length $h_m$, an observation and not a result.

\paperfigure{E1}{\figcapEOne}{fig:E1}

\subsection{The systole: an open question}

Pinching a closed geodesic produces small eigenvalues, but this does not show that the systole bound is needed. Fix $k\in\{3,4\}$, $b>0$ and $a$ with $\sinh a\sinh b=\cos\frac\pi k$; in $\Hb$ take orthogonal lines $\beta,\beta'$ through $O$, and let $Q_{k,b}$ be the quadrilateral bounded by the two lines perpendicular to $\beta$ at distance $b$ from $O$ and the two perpendicular to $\beta'$ at distance $a$. In the hyperboloid, with $O=(1,0,0)$ and $\beta,\beta'$ in the planes $x_2=0$, $x_1=0$, the unit normals $(\sinh b,\cosh b,0)$ and $(\sinh a,0,\cosh a)$ of adjacent sides have inner product $-\sinh a\sinh b$, of modulus less than $1$, so these sides meet at an angle with cosine $\sinh a\sinh b$ or its negative (Lemma~\ref{lem:H}, proof of (H3)); the quarter of $Q_{k,b}$ cut out by $\beta,\beta'$ has three right angles and angle sum less than $2\pi$, so its fourth angle is acute, equal to $\frac\pi k$. Let $\Orb_{k,b}$ be the double of $Q_{k,b}$, a sphere with four cone points of order $k$.

\begin{proposition}[Pinching]\label{prop:N2}
$\Orb_{k,b}$ contains an embedded collar $\{|\rho|<a\}\times\R/4b\Z$, with metric $d\rho^2+\cosh^2\rho\,ds^2$, about a separating closed geodesic of length $4b$; so $\ell(\Orb_{k,b})\le4b$, $\lambda_1(\Orb_{k,b})\le\sinh a/((a^2+2)\sinh a-2a\cosh a)$ and $\lambda_j(\Orb_{k,b})\le\frac14+\frac1{4\cosh^2(a/2)}+\frac{4\pi^2(j+1)^2}{a^2}$. As $b\to0$, $\lambda_1\to0$ and $\limsup\lambda_j\le\frac14$.
\end{proposition}

\begin{proof}
In Fermi coordinates $(\rho,s)$ about $\beta$ the sides perpendicular to $\beta$ are $s=\pm b$, and the other two sides are at distance at least $a$ from $\beta$, their common perpendiculars with $\beta$ having length $a$; so $\{|\rho|<a,|s|\le b\}$ lies in $Q_{k,b}$, and its double is the collar, whose core, the double of $\beta\cap Q_{k,b}$, separates the two pairs of cone points. The reflection in $\beta$ induces an isometry exchanging the two sides. For $\lambda_1$ take $f=\rho/a$ on the collar and $\pm1$ outside, odd under that isometry, so $\int f=0$, and compute the Rayleigh quotient; for $\lambda_j$ use \eqref{eq:rayleigh} with $w=\cosh\rho$ and $j+1$ sine arches on disjoint intervals of length $a/(2(j+1))$ in $[\frac a2,a)$, where $q\le\frac14+\frac1{4\cosh^2(a/2)}$.
\end{proof}

Since $Q_{k,b}$ is symmetric under $a\leftrightarrow b$, $\lambda_1(\Orb_{k,b})\to0$ also as $b\to\infty$, and the intermediate value theorem gives $b,b'$ in a compact interval with $\lambda_1(\Orb_{3,b})=\lambda_1(\Orb_{4,b'})$: two eigenvalues fail to decide the signature already at a fixed systole, which only says that $N=2$ is too small.

\begin{problem}[Is a systole bound needed?]\label{prob:systole}
Fix $A$ and $M$. Can $N$ and $\delta$ in Theorem~\ref{thm:E} be chosen independently of $\varepsilon$?
\end{problem}

With $A$ and $M$ fixed only finitely many signatures occur, so the pigeonhole argument of Proposition~\ref{prop:N1} is not available, and the number of eigenvalues that pinching drives to $0$ is bounded in terms of $A$. A negative answer needs two degenerating families of different signatures whose low spectra converge to the same limit: the upper bounds above give $\limsup\lambda_j\le\frac14$, and the missing input is a matching lower bound $\lambda_j\ge\frac14-o(1)$ for $2\le j<N$ along both families. Such a bound might come from finite manifold covers of bounded degree, whose spectra contain those of the orbifolds, together with the spectral theory of degenerating hyperbolic surfaces; we have not carried this out. The $N$ of Theorem~\ref{thm:E} grows like $\varepsilon^{-3}\log\frac1\varepsilon$ as $\varepsilon\to0$ once $t_3$ binds.

\backmatter

\bmhead{Acknowledgements}
The colour schemes of the figures are Scientific Colour Maps \cite{crameri2023,crameri2020}. The eigenvalues were computed with NGSolve \cite{ngsolve,schoberl1997} and ARPACK \cite{arpack1998}.

\section*{Statements and Declarations}

\bmhead{Funding}
No funding was received for this work.

\bmhead{Competing interests}
The authors have no competing interests to declare that are relevant to the content of this article.

\bmhead{Ethics approval and consent}
Not applicable.

\bmhead{Data and code availability}
All code and data are in the public repository \url{https://github.com/Ali-M658/Arithmetic-verification}. One command, \texttt{bash code/run\_all.sh}, reruns every check (with \texttt{--quick} for the exact-arithmetic suite and \texttt{--full} for the long computations, including the eigenvalue computations of Figure~\ref{fig:E1}), and \texttt{DATA-MANIFEST.md} lists every data file with its generating script and checksum. Every constant of Tables~\ref{tab:D}--\ref{tab:practice} is recomputed there at $40$ digits. An archived copy will be deposited at Zenodo: \textbf{[PLACEHOLDER: Zenodo DOI, to be minted at submission.]}

\bmhead{Author contributions}
\textbf{[PLACEHOLDER: author contribution statement, to be supplied by the authors before submission.]}

\bmhead{Use of AI tools}
\textbf{[PLACEHOLDER: statement on the use of AI tools in preparing this manuscript, to be written by the authors. Springer Nature's guidelines ask for any such use, beyond AI-assisted copy editing, to be documented in the Methods section or, where there is none, in a suitable alternative part of the manuscript.]}

\begin{appendices}

\section{The trace formula for the heat function}\label{app:trace}

Dryden and Strohmaier state Theorem~\ref{thm:trace} for $h=\hat g$ with $g\in C_c^\infty(\R)$ even \cite[eq.~(1)]{drydenstrohmaier2009}. The following standard approximation extends it to the heat function.

\begin{lemma}[The heat function is admissible]\label{lem:admissible}
Fix an even $\omega\in C_c^\infty(\R)$, $0\le\omega\le1$, $\omega=1$ on $[-1,1]$, and put $h_\varrho=\hat g_\varrho$ with $g_\varrho=g_t\,\omega(\cdot/\varrho)$. As $\varrho\to\infty$, every term of the trace formula for $h_\varrho$ tends to the corresponding term for $h_t$; hence Theorem~\ref{thm:trace} holds.
\end{lemma}

\begin{proof}
\emph{Counting eigenvalues.} Let $\lambda_j=\frac14+r_j^2$ and take $\psi\in C_c^\infty(\R)$ even, nonnegative, $\int\psi=1$, supported in $[-\epsilon,\epsilon]$ with $2\epsilon<\ell(\Orb)$, $\epsilon\le1$; then $|\hat\psi|\le1$ and $\hat\psi(r)\ge\cos(\epsilon r)\ge\frac12$ for $|r|\le1$. The function $\hat\psi(r-T)^2+\hat\psi(r+T)^2$ is $\hat g$ for $g=2\cos(Tu)(\psi*\psi)(u)$, supported in $[-2\epsilon,2\epsilon]$, so the hyperbolic side vanishes for it; its identity term is $O(T+1)$, each elliptic term $O(1)$, and the finitely many imaginary $r_j$ contribute $O(1)$. As it is at least $\frac14$ on $|r-T|\le1$, $\#\{j:|r_j-T|\le1\}=O(T+1)$, so $\#\{\lambda_j\le x\}=O(x)$.
\emph{Convergence.} On the hyperbolic side $0\le g_\varrho\le g_t$, with equality on $[0,\varrho]$, and Lemma~\ref{lem:counting} gives dominated convergence. With $f_\varrho=g_t(1-\omega(\cdot/\varrho))$, $h_t-h_\varrho=\hat f_\varrho$ and $\|f_\varrho^{(k)}\|_{L^1}\to0$ for $k\le3$, so $|\hat f_\varrho(r)|\le\epsilon_\varrho\min(1,|r|^{-3})$ with $\epsilon_\varrho\to0$; this handles the identity and elliptic terms and, as $\sum_j\min(1,|r_j|^{-3})<\infty$, the real $r_j$; for imaginary $r_j$, $|\hat f_\varrho(r_j)|\le\int|f_\varrho|e^{|u|/2}\to0$.
\end{proof}

\end{appendices}

\bibliography{references}

\end{document}
